% 2. 编译器的安装
\section{编译器的安装}

\begin{frame}
    \frametitle{编译器的安装}
    安装方式按平台分三类：
    \begin{itemize}
        \item \alert{Linux}：\code{gcc}/\code{g++}
        \item \alert{Windows}：通过WSL使用\code{gcc}/\code{g++}
        \item \alert{Mac}：自带\code{clang}/\code{clang++}
    \end{itemize}
\end{frame}

\begin{frame}
    \frametitle{Linux（以Ubuntu为例）}
    一行命令即可安装编译所需的一整套工具（\code{gcc}、\code{g++}、\code{make}等）：
    \begin{Code}*
        \lstinputlisting{code/setup/install-build-essential.sh}
    \end{Code}

    在Ubuntu24.04默认安装的是gcc-13，26.04默认安装的是gcc-15。

    每个新版本的GCC都会增加对更新标准的支持：
    \begin{itemize}
        \item gcc-9：对C++17支持比较完善；
        \item gcc-11：改进了C++20的支持；
        \item gcc-13/14：开始逐步实现C++23的特性。
    \end{itemize}
    同时编译器优化更强，错误检查与诊断也更好。
\end{frame}

\begin{frame}
    \frametitle{Windows（WSL）}
    WSL（Windows Subsystem for Linux，Windows的Linux子系统）是微软官方提供的功能，本质上是一个在Windows上运行的\alert{只有命令行环境的Linux虚拟机}。

    换句话说：在Windows里，你可以直接打开一个Linux终端，运行\code{apt}、\code{bash}、\code{gcc}、\code{python}等工具，就像真的在Linux下一样。

    \alert{从这学期开始，课程演示全部在Linux环境下进行}：Windows用户请全程使用WSL，Mac用户直接使用自带终端。
\end{frame}

\begin{frame}
    \frametitle{WSL安装}
    WSL的安装需要在\alert{正确的网络环境}下，在管理员权限的PowerShell中执行：
    \begin{Code}*
        \lstinputlisting{code/setup/wsl-install.sh}
    \end{Code}
    会自动安装WSL和默认的Ubuntu，装完\alert{重启电脑}生效。

    如果网络环境不正确，可以先去GitHub的\code{microsoft/WSL} releases下载最新\alert{x86}架构的\code{.msi}手动装上WSL本体\footnote{次选方案：如果实在没有正确的网络环境，可以试试这个办法。}，再执行\code{wsl --install}安装Ubuntu；若安装卡在0\%，可加参数\code{--web-download}。

    装完Ubuntu后首次打开，会让你创建一个用户名和密码（注意密码输入时屏幕不显示字符，是正常的）。
\end{frame}

\begin{frame}
    \frametitle{WSL安装的版本选择}
    推荐使用Ubuntu 26.04或Ubuntu 24.04，两种版本本课程都支持：
    \begin{itemize}
        \item \code{wsl --install}：Ubuntu 26.04（默认gcc-15）；
        \item \code{wsl --install Ubuntu-24.04}：Ubuntu 24.04（默认gcc-13）。
    \end{itemize}

    装好后用\code{cat /etc/os-release}确认版本，后面换源时选择对应的脚本。
\end{frame}

\begin{frame}
    \frametitle{双系统：更好的选择}
    WSL只有命令行环境，如果你希望体验完整的Linux桌面，可以考虑\alert{双系统}：
    \begin{itemize}
        \item 在磁盘上划出一块空间，在保留Windows的情况下安装完整的Linux系统；
        \item 开机时可以选择启动哪个系统；
        \item 推荐Ubuntu 26.04或24.04，也可以使用其他发行版。
    \end{itemize}

    我们真诚鼓励所有学生尝试双系统方案，体验在一个完整的Linux桌面系统上工作。
\end{frame}

\begin{frame}
    \frametitle{WSL环境配置：换源}
    首次进入WSL后，先把软件源换到国内镜像，否则\code{apt}下载很慢。

    先执行\code{cat /etc/os-release}确认Ubuntu版本（看\texttt{VERSION\_CODENAME}），然后按版本选择换源脚本：
    \begin{itemize}
        \item Ubuntu24.04（noble）：\texttt{magic\_code\_2404.md}
        \item Ubuntu26.04（resolute）：\texttt{magic\_code\_2604.md}
    \end{itemize}

    \begin{vspquotered}{版本必须匹配}
        24.04的机器用26.04的源（或反过来）会导致仓库里所有包都比系统旧，\code{apt install}会报依赖错误。
    \end{vspquotered}
\end{frame}

\begin{frame}
    \frametitle{WSL环境配置：安装编译工具}
    \begin{vspquote}{WSL与Linux配置编译环境}
        按顺序执行，\code{build-essential}会一次装齐\code{gcc}、\code{g++}、\code{make}等工具：
        \begin{Code}*
            \lstinputlisting{code/setup/apt-setup.sh}
        \end{Code}
    \end{vspquote}
\end{frame}

\begin{frame}
    \frametitle{Mac环境配置}
    Mac不用额外配置，系统自带\code{clang}编译器。用下面的命令检查版本即可：
    \begin{Code}*
        \lstinputlisting{code/setup/clang-version.sh}
    \end{Code}
\end{frame}
